\section{Fundamental Equations of I}
\label{sec:fundamental_equations}

This section presents the fundamental physical equations governing the behavior of the indeterminate I as a physical entity.

\subsection{Mass-Energy Relation}
\label{subsec:mass_energy}

In analogy to standard particles, the mass-energy equivalence for I is:

\begin{equation}
\label{eq:i_mass_energy}
E_I = m_I c^2
\end{equation}

where:
\begin{itemize}
    \item $m_I$ is the rest mass of the I particle
    \item $c$ is the speed of light in vacuum ($c = 299{,}792{,}458$ m/s)
\end{itemize}

This relation indicates that if I carries a finite mass, it can in principle be created or annihilated by processes that furnish sufficient energy.

\subsubsection{Particle Production Threshold}

The threshold energy for I-particle pair production:

\begin{equation}
\label{eq:i_production_threshold}
E_{\text{threshold}} = 2m_I c^2
\end{equation}

In particle collisions with center-of-mass energy $E_{\text{CM}}$:

\begin{equation}
\label{eq:i_production_condition}
E_{\text{CM}} \geq 2m_I c^2 + E_{\text{binding}}
\end{equation}

where $E_{\text{binding}}$ accounts for any binding energy in I-composite states.

\subsubsection{I-Mass Bounds}

From cosmological observations (dark matter constraints):

\begin{equation}
\label{eq:i_mass_lower_bound}
m_I \gtrsim 1 \text{ keV}/c^2
\end{equation}

From collider searches:

\begin{equation}
\label{eq:i_mass_upper_bound}
m_I \lesssim 10^3 \text{ GeV}/c^2 \quad \text{(model-dependent)}
\end{equation}

\subsubsection{Relativistic Generalization}

For moving I-particles:

\begin{equation}
\label{eq:i_relativistic_energy}
E_I = \gamma m_I c^2 = \frac{m_I c^2}{\sqrt{1 - v^2/c^2}}
\end{equation}

The momentum-energy relation:

\begin{equation}
\label{eq:i_momentum_energy}
E_I^2 = (p_I c)^2 + (m_I c^2)^2
\end{equation}

\subsection{Gravitational Interaction}
\label{subsec:gravitational_interaction}

If I has a non-negligible mass, it experiences a classical gravitational force when near a massive body of mass $M$:

\begin{equation}
\label{eq:gravitational_force}
F_G = G \frac{m_I M}{r^2}
\end{equation}

where:
\begin{itemize}
    \item $G$ is the gravitational constant ($G = 6.674 \times 10^{-11}$ m$^3$ kg$^{-1}$ s$^{-2}$)
    \item $r$ is the separation between the center of mass $M$ and the I particle
    \item The force is attractive (directed toward the mass $M$)
\end{itemize}

\subsubsection{Gravitational Potential Energy}

The potential energy of an I-particle in a gravitational field:

\begin{equation}
\label{eq:gravitational_potential}
U_G(r) = -G\frac{m_I M}{r}
\end{equation}

The total energy:

\begin{equation}
\label{eq:total_energy}
E_{\text{total}} = \frac{1}{2}m_I v^2 - G\frac{m_I M}{r} = \text{constant}
\end{equation}

\subsubsection{Orbital Dynamics}

For I-particles in orbit around a central mass:

\begin{equation}
\label{eq:orbital_velocity}
v_{\text{orbit}} = \sqrt{\frac{GM}{r}}
\end{equation}

The orbital period:

\begin{equation}
\label{eq:orbital_period}
T = 2\pi\sqrt{\frac{r^3}{GM}}
\end{equation}

\subsubsection{Cosmic Applications}

In cosmic contexts, gravitational attraction could influence:

\begin{enumerate}
    \item Orbital motion of I-particles in galactic halos
    \item Clustering properties in large-scale structure
    \item Accretion onto compact objects (black holes, neutron stars)
\end{enumerate}

The virial theorem for I-particle systems:

\begin{equation}
\label{eq:virial_theorem}
2\langle K_I \rangle = -\langle U_G \rangle
\end{equation}

where $\langle K_I \rangle$ is the average kinetic energy and $\langle U_G \rangle$ is the average gravitational potential energy.

\subsection{Energy Density in a Dark Halo}
\label{subsec:energy_density}

For a region in which numerous I particles congregate, the effective mass density can be approximated by:

\begin{equation}
\label{eq:i_density}
\rho_I = \frac{\sum m_I}{V}
\end{equation}

where:
\begin{itemize}
    \item The summation $\sum m_I$ runs over all particles within the volume $V$
    \item $\rho_I$ reflects a dark matter-like density distribution if I dominates the local mass content
\end{itemize}

\subsubsection{Spatial Distribution}

For a spherically symmetric halo:

\begin{equation}
\label{eq:halo_profile}
\rho_I(r) = \frac{\rho_0}{\left(r/r_s\right)^\gamma\left(1 + r/r_s\right)^{3-\gamma}}
\end{equation}

Common profiles:
\begin{itemize}
    \item NFW profile: $\gamma = 1$
    \begin{equation}
    \label{eq:nfw_profile}
    \rho_{\text{NFW}}(r) = \frac{\rho_0}{(r/r_s)(1 + r/r_s)^2}
    \end{equation}
    
    \item Isothermal sphere: $\rho(r) \propto r^{-2}$
    \begin{equation}
    \label{eq:isothermal_profile}
    \rho_{\text{iso}}(r) = \frac{v_{\text{rot}}^2}{4\pi G r^2}
    \end{equation}
\end{itemize}

\subsubsection{Total Mass}

The total I-mass within radius $R$:

\begin{equation}
\label{eq:enclosed_mass}
M_I(R) = 4\pi \int_0^R \rho_I(r) r^2 \, dr
\end{equation}

For an NFW profile:

\begin{equation}
\label{eq:nfw_mass}
M_I(R) = 4\pi\rho_0 r_s^3 \left[\ln\left(1 + \frac{R}{r_s}\right) - \frac{R/r_s}{1 + R/r_s}\right]
\end{equation}

\subsubsection{Velocity Dispersion}

The velocity dispersion of I-particles:

\begin{equation}
\label{eq:velocity_dispersion}
\sigma_v^2 = \frac{1}{m_I\rho_I(r)} \int_r^\infty \rho_I(r') \frac{GM_I(r')}{r'^2} dr'
\end{equation}

For an isothermal distribution:

\begin{equation}
\label{eq:isothermal_dispersion}
\sigma_v = \text{constant} = v_{\text{rot}}
\end{equation}

\subsection{Interaction Cross-Sections}

\subsubsection{I-I Scattering}

The differential cross-section for I-particle scattering:

\begin{equation}
\label{eq:i_i_scattering}
\frac{d\sigma}{d\Omega} = \left|\frac{f(\theta)}{2ik}\right|^2
\end{equation}

where $f(\theta)$ is the scattering amplitude.

For weak interactions:

\begin{equation}
\label{eq:weak_scattering}
\sigma_{I-I} = \frac{G_F^2 m_I^2}{\pi}
\end{equation}

where $G_F$ is the Fermi constant.

\subsubsection{I-Baryon Interactions}

Direct detection constraints:

\begin{equation}
\label{eq:direct_detection}
\sigma_{I-N} \lesssim 10^{-45} \text{ cm}^2 \quad \text{(for } m_I \sim 100 \text{ GeV)}
\end{equation}

Spin-independent interaction:

\begin{equation}
\label{eq:spin_independent}
\mathcal{L}_{I-N} = \frac{g_{IN}}{m_N^2} \bar{I}I \bar{N}N
\end{equation}

\subsection{Thermodynamics of I-Systems}

\subsubsection{Maxwell-Boltzmann Distribution}

For non-relativistic I-particles in thermal equilibrium:

\begin{equation}
\label{eq:maxwell_boltzmann}
f(v) = \left(\frac{m_I}{2\pi k_B T}\right)^{3/2} \exp\left(-\frac{m_I v^2}{2k_B T}\right)
\end{equation}

The most probable speed:

\begin{equation}
\label{eq:most_probable_speed}
v_{\text{mp}} = \sqrt{\frac{2k_B T}{m_I}}
\end{equation}

\subsubsection{Equation of State}

For an ideal gas of I-particles:

\begin{equation}
\label{eq:ideal_gas}
P = \frac{\rho_I k_B T}{m_I} = \frac{n_I k_B T}{\text{NA}}
\end{equation}

The sound speed:

\begin{equation}
\label{eq:sound_speed}
c_s = \sqrt{\frac{\gamma k_B T}{m_I}}
\end{equation}

where $\gamma = C_P/C_V$ is the adiabatic index.

\subsubsection{Jeans Instability}

The Jeans mass for I-particle collapse:

\begin{equation}
\label{eq:jeans_mass}
M_J = \left(\frac{5k_B T}{Gm_I}\right)^{3/2} \left(\frac{3}{4\pi\rho_I}\right)^{1/2}
\end{equation}

The Jeans length:

\begin{equation}
\label{eq:jeans_length}
\lambda_J = c_s \sqrt{\frac{\pi}{G\rho_I}}
\end{equation}

\subsection{Quantum Corrections}

\subsubsection{Fermi-Dirac Statistics}

If I-particles are fermions:

\begin{equation}
\label{eq:fermi_dirac}
f(E) = \frac{1}{e^{(E-\mu)/k_B T} + 1}
\end{equation}

The Fermi energy:

\begin{equation}
\label{eq:fermi_energy}
E_F = \frac{\hbar^2}{2m_I}(3\pi^2 n_I)^{2/3}
\end{equation}

\subsubsection{Bose-Einstein Statistics}

If I-particles are bosons:

\begin{equation}
\label{eq:bose_einstein}
f(E) = \frac{1}{e^{(E-\mu)/k_B T} - 1}
\end{equation}

Below the critical temperature:

\begin{equation}
\label{eq:bec_critical_temperature}
T_c = \frac{2\pi\hbar^2}{m_I k_B}\left(\frac{n_I}{\zeta(3/2)}\right)^{2/3}
\end{equation}

I-particles condense into the ground state.

\subsection{Field-Theoretic Description}

\subsubsection{Klein-Gordon Equation}

For scalar I-fields:

\begin{equation}
\label{eq:klein_gordon}
\left(\Box + \frac{m_I^2 c^2}{\hbar^2}\right)\phi_I = 0
\end{equation}

where $\Box = \frac{1}{c^2}\frac{\partial^2}{\partial t^2} - \nabla^2$ is the d'Alembertian operator.

\subsubsection{Dirac Equation}

For fermionic I-fields:

\begin{equation}
\label{eq:dirac}
\left(i\hbar\gamma^\mu\partial_\mu - m_I c\right)\psi_I = 0
\end{equation}

where $\gamma^\mu$ are the Dirac matrices.

\subsubsection{Lagrangian Density}

The complete I-field Lagrangian:

\begin{equation}
\label{eq:i_field_lagrangian}
\mathcal{L}_I = \frac{1}{2}\partial_\mu\phi_I\partial^\mu\phi_I - \frac{1}{2}m_I^2\phi_I^2 - \frac{\lambda}{4!}\phi_I^4 - V_{\text{int}}(\phi_I, \phi_{\text{SM}})
\end{equation}

where $V_{\text{int}}$ describes interactions with Standard Model fields.
